\label{exp:cea02}

\subsubsection*{Постановка}
\textbf{Цель.} Статистика head-on T2: ось mod~$\pi$, интервал Стьюдента для $\bar\varepsilon_\pi$; отделить $\tau_M$ от $\tau_{\mathrm{PDG}}$.

\textbf{Критерий.} PASS \texttt{Annihilation\_T\_stats}; интервал $\bar\theta\pm t_{n-1,0.975}\,\mathrm{SE}$ (\cref{sec:student}).

\subsubsection*{Теоретическая часть}
Ансамбль вращений инъекции; $\tau_M=n_{\mathrm{ticks}} hT$ --- часы столкновения, не para-Ps (\cref{ch:measurement-processing}).

\subsubsection*{Ход эксперимента}
\texttt{python verify\_principles.py --suite sm}; $n=12$ в реестре.

\begin{table}[htbp]
\centering
\small
\caption{CE-A-02}
\label{tab:exp-ce-a-02}
\begin{tabular}{@{}ll@{}}
\toprule
Наблюдаемое & Значение \\
\midrule
$\bar\varepsilon_\pi$ [rad] & $6{,}37\times 10^{-2}$ \\
SEM [rad] & $1{,}03\times 10^{-2}$ \\
$\tau_M/\tau_{\mathrm{PDG}}$ & $4{,}88\times 10^{-33}$ \\
back-to-back & True \\
\bottomrule
\end{tabular}
\end{table}

\subsubsection*{Анализ, расчёт погрешности и выводы}
$t_{11,0.975}=2{,}201$; полуинтервал $\approx 2{,}26\times 10^{-2}\,\mathrm{rad}$.
Тип~A; не спектроскопия PDG.

\textbf{Статус:} закрыт.
